% Copyright (c) 2008-2009 solvethis
% Copyright (c) 2010-2016,2018-2019,2022 Casper Ti. Vector
% Copyright (c) 2021 Kurapica
% Copyright (c) 2022 iofu728
% Overleaf version.

%*********************************************************************
% iofu728-pkuthss: 北京大学研究生学位论文模板
% 2022/07/16 v1.1.0
%
% 重要提示：
%   1. 当前overleaf版符合2022研究生学位论文要求，可通过图书馆审核
%   2. 当前版本基于pkuthss v1.9.2
%   3. 请使用UTF-8编码，XeLaTeX方式编译
%   4. 请仔细阅读用户文档
%   5. 修改、使用、发布本文档类请务必遵循LaTeX Project Public License和知识共享4.0
%   6. 如有疑问github/iofu728/pkuthss上提问或联系作者@iofu728
%*********************************************************************

\PassOptionsToPackage{shortlabels}{enumitem}
\documentclass[fontset=fandol,ugly,language=chinese]{pkuthss}
  % 学位论文模式  ugly    (默认打开，请保留)
  % 盲审模式     blind   (默认关闭)
  % 语言        language (默认chinese | english)
  % 字体库       fontset
  %   auto | windows | windows@overleaf | mac | fandol | ubuntu | none
  % windows*, mac为商业字体，如需使用请遵循相应版权协议(默认下overleaf中不可用)
  % fandol与windows效果相近，但字符库偏少，推荐使用(默认);
  % ubuntu字体效果偏差较大; 设为none时需自行配置字体集;

\usepackage[backend=bibtex,style=gb7714-2015]{biblatex}
  % 参考文献遵循GB/T 7714-2015标准，使用biblatex-gb7714-2015 宏包。
  % 此处使用顺序编码制，如使用著者-出版年制则更改为b7714-2015ay。
  % 注：MiKTeX biber 存在 Windows 路径兼容问题，改用 bibtex 后端。

% 工程硕士论文常用宏包
\usepackage{booktabs,multirow,longtable,makecell}
\usepackage{amsmath}
\usepackage{tikz}
\usetikzlibrary{arrows.meta, shapes.geometric, shapes.multipart, positioning, fit, backgrounds, calc, chains}
\renewcommand{\v}[1]{\boldsymbol{#1}}
\newcommand\pkg[1]{\textsf{#1}}

% 参考文献边距字体
\setlength{\bibitemsep}{3bp}
\renewcommand*{\bibfont}{\zihao{5}\linespread{1.27}\selectfont}

\pkuthssinfo{
    ctitle = {基于权重驻留脉动阵列的轻量级卷积神经网络FPGA加速器设计},
    etitle = {Design of a Lightweight CNN FPGA Accelerator Based on Weight-Stationary Systolic Array},
}

\addbibresource{ref.bib}

% 重新定义页眉：用章节标题替代默认的"北京大学硕士学位论文"
\fancypagestyle{plain}{
    \fancyhf{}
    \renewcommand*{\headrulewidth}{0.75bp}
    \fancyfoot[C]{\zihao{5}\normalfont{\thepage}}
    \fancyhead[C]{\zihao{5}\normalfont\leftmark}
}

\begin{document}
    \frontmatter
    \pagestyle{empty}
    \setcounter{page}{0}
    \pagenumbering{Roman}

    \chapter*{摘要}
\markboth{摘要}{}

边缘端目标检测要求在有限功耗和有限片上资源下完成较低延迟的卷积神经网络推理。FPGA 与 MPSoC 平台具有可定制数据通路、丰富片上存储和灵活软硬件协同的优势，适合部署轻量级目标检测网络。然而，在真实网络端到端运行中，加速器性能并不只取决于乘加阵列规模；DMA 启动、输入特征图重排、权重加载、部分和反馈、多轮 K pass 同步以及软硬件调度空泡都会显著降低阵列有效利用率。

本文面向 Xilinx Kria KV260 / XCK26 平台，设计并实现了一种用于 YOLOv3-tiny 单尺度链式推理的 INT8 卷积加速器。硬件采用参数化 systolic array 作为卷积核心，支持 3x3 卷积、原生 1x1 卷积、INT32 部分和累加、逐通道重量化、查表激活、池化以及 AXI-Lite/AXI-Stream 系统接口。软件侧基于 Vitis bare-metal runtime 负责层描述、空间 tile 调度、DMA 启动、量化参数配置、输出重排和 YOLO decode，从而形成从 DDR 图像输入到 Conv0--Conv9 输出解析的完整 PS/PL 协同链路。

围绕真实 YOLO 层的非计算开销，本文进一步提出并实现了多级数据流与调度优化。首先，通过 batch stream 和预打包权重布局减少每个 tile 内的 DMA 启动与运行时重排开销；其次，引入 raw-HWC tile cache，使软件能够传输自然 HWC tile，由硬件在片上完成 1x1/3x3 向量重放，并在后端层使用 full-tile 调度降低空间 tile 启动开销；再次，针对多 K pass 卷积中的 partial-PSUM 反馈瓶颈，实现了流水化 drain、early drain、continuous PSUM collector 和 pass-level prefetch，使权重准备、IFM replay、部分和反馈与当前 pass 计算尽可能重叠。本文同时设计了阶段、子阶段、pass 级和列级性能计数器，用于解释瓶颈迁移和验证优化有效性。

实验结果表明，所设计系统在 KV260 上完成 Conv0--Conv9 batch-chain bit-exact 验证，并在两张 DDR 动态图像上得到稳定的 YOLO decode 结果。最终推荐实现采用 Vivado/Vitis 2022.2、100 MHz 时钟、ROWS=18、COLS=8、COUT\_TILE=16 的配置，在 D:/MPSoC/b\_kprefetch\_22 构建中完成布局布线，资源占用为 84480 个 CLB LUT、53260 个 CLB Register、63 个 BRAM Tile、24 个 URAM 和 184 个 DSP。两张图像的端到端 DDR demo 延迟约为 288 ms。该结果并不定位为实时视频性能领先，而是证明在资源受限 MPSoC 上，围绕数据流、部分和反馈和 pass 级调度的系统优化能够将早期约 1.18 s 的链式推理延迟降低到约 0.29 s，并给出可解释的瓶颈分析路径。

\medskip
\noindent\textbf{关键词：}YOLOv3-tiny，MPSoC，FPGA 加速器，systolic array，数据流优化，partial-PSUM，INT8 量化

    \cleardoublepage
    \tableofcontents\markboth{目录}{}
    \cleardoublepage
    \listoftables\markboth{表格索引}{}
    \cleardoublepage
    \listoffigures\markboth{插图索引}{}
    \cleardoublepage
    \begin{denotation}[6.5cm]

\item[CNN] 卷积神经网络 (Convolutional Neural Network)
\item[IFM] 输入特征图 (Input Feature Map)
\item[OFM] 输出特征图 (Output Feature Map)
\item[PSUM] 部分和 (Partial Sum)
\item[DSP] 数字信号处理单元 (Digital Signal Processing slice)
\item[BRAM] 块随机存取存储器 (Block RAM)
\item[LUT] 查找表 (Look-Up Table)
\item[PE] 处理单元 (Processing Element)
\item[FSM] 有限状态机 (Finite State Machine)
\item[AXI] 高级可扩展接口 (Advanced eXtensible Interface)
\item[DMA] 直接存储器访问 (Direct Memory Access)
\item[$K_{\text{total}}$] 展开后输入通道总维度，$K_{\text{total}} = C_{in} \times K_h \times K_w$
\item[$K_{\text{TILE}}$] 单次 K 分块处理的输入 lane 数
\item[$C_{\text{out}}$] 输出通道数
\item[$\text{COUT}_{\text{TILE}}$] 单次计算覆盖的输出通道数
\item[$H, W, C$] 特征图高度、宽度、通道数
\item[$K_h, K_w$] 卷积核高度、宽度
\item[$oy, ox$] 输出特征图空间坐标
\item[GEMM] 通用矩阵乘法 (GEneral Matrix Multiply)
\item[FPGA] 现场可编程门阵列 (Field Programmable Gate Array)
\item[XCK26] Xilinx Kria K26 核心 FPGA 器件型号
\item[WNS] 最差负时序裕量 (Worst Negative Slack)
\item[HWC] 高度-宽度-通道 (Height-Width-Channel) 内存布局

\end{denotation}


    \mainmatter
    \pagestyle{plain}
    \chapter{绪论}
\label{chap:introduction}

\section{研究背景与意义}

卷积神经网络（Convolutional Neural Network, CNN）已经成为目标检测、图像分类和语义理解等视觉任务的核心计算模型。YOLO 系列检测器将候选框生成和类别预测整合为单阶段推理流程，在边缘端实时感知任务中具有较好的工程适配性~\cite{redmon2016yolo,redmon2018yolov3}。其中 YOLOv3-tiny 通过减少层数和通道规模降低计算量，适合在嵌入式平台上部署。

边缘端部署面对的约束与数据中心不同。通用 CPU 的吞吐有限，嵌入式 GPU 或专用 NPU 虽然性能较高，但在功耗、可定制性、接口灵活性和可验证性方面并不总是满足定制系统需求。FPGA 与 MPSoC 平台同时包含处理系统（Processing System, PS）和可编程逻辑（Programmable Logic, PL），能够让软件负责图像输入、调度和后处理，让硬件承担规则密集的卷积计算。因此，基于 FPGA 的 CNN 加速器仍然是边缘智能系统中的重要研究方向~\cite{qiu2016going,guo2019angeleye}。

另一方面，真实检测网络的层结构具有明显异构性。以本文采用的 YOLOv3-tiny Conv0--Conv9 链路为例，前端层空间尺寸较大且常伴随池化，后端层空间尺寸较小但通道数显著增加；网络同时包含 3x3 卷积和 1x1 检测头，不同层对 activation、pooling、K pass 数量和输出通道分块的需求也不相同。若硬件仅针对单一卷积核或固定后处理路径设计，则 1x1 层可能被低效映射为稀疏 3x3，池化和激活需要频繁回退到软件，或者每类层都需要独立硬件路径，最终降低资源利用率和系统可扩展性。因此，本文首先需要构建一套层自适应的卷积执行基础，使同一套脉动阵列能够通过运行时描述符适配不同层形状和后处理组合。

然而，仅具备层自适应能力仍不足以保证端到端性能。真实 CNN 网络的加速并不等同于构造一个高峰值 MAC 阵列；每一层还要经过输入通道展开、输出通道分块、空间 tile 分块和多轮 K pass 累加。若数据搬运、部分和反馈或 pass 边界调度处理不当，硬件阵列会在等待权重、等待 IFM replay、等待 partial-PSUM drain 或等待 DMA 完成时产生大量空泡。本文关注的核心问题正是：\textbf{如何在资源受限的 MPSoC 上，构建层自适应的卷积加速基础，并面向真实 YOLOv3-tiny 链式推理降低非计算开销，使优化效果能够被板级实验和性能计数器解释。}

\section{国内外研究现状}

CNN 加速器的研究大致经历了从算子级流水线到空间阵列架构的发展。Kung 提出的 systolic architecture 通过规则 PE 网络在相邻单元之间传递数据，为矩阵乘和卷积计算提供了高复用的硬件组织方式~\cite{kung1982systolic}。Google TPU 使用大规模 systolic array 获得高吞吐推理性能~\cite{jouppi2017tpu}。Eyeriss 系列则从能量效率角度系统讨论了 row-stationary 数据流和层级存储复用~\cite{chen2016eyeriss,chen2019eyerissv2}。Sze 等人的综述对 DNN 加速器的数据流、存储层次和能效权衡进行了系统总结~\cite{sze2017efficient}。

在 FPGA 方向，Qiu 等人较早展示了在嵌入式 FPGA 上部署 CNN 的完整流程~\cite{qiu2016going}；Angel-Eye 进一步提出从模型到 FPGA 映射的设计流程~\cite{guo2019angeleye}；S2A 和 RTL compilation 相关工作关注从高层描述生成 systolic array 或模块化 RTL~\cite{wei2017s2a,ma2018rtlcompiler}。针对 YOLOv3-tiny 的轻量化 FPGA 加速器也已有报道~\cite{ding2022tinyyolo}。这些工作说明 FPGA 对 CNN 推理具有可行性，但许多研究更强调模型映射、阵列生成或单层性能，对真实链式推理中的 DMA 交互、IFM 打包、partial-PSUM 反馈和 pass 级调度开销讨论较少。

本文以可复现的 MPSoC 端到端系统为载体，研究真实 YOLO 链式推理中的数据流和调度瓶颈。相比仅报告理论吞吐或单层仿真的工作，本文更强调硬件 RTL、Vitis runtime、外部 golden、板级 batch-chain 验证、DDR 图像输入和性能计数器之间的一致闭环。

图~\ref{fig:workflow_overview} 给出本文的系统框架与优化路线。论文首先从真实 CNN 链式推理中的数据搬运和调度开销出发，构建 PS/PL 协同的 LASA 卷积执行引擎；随后以 KCS 三维分块模型统一层映射，并通过 BSD、OCRR 和 OPF-P 分别优化 PS/PL 交互、IFM 数据排布和 K-pass 边界空泡；最终通过 golden 验证、DDR image demo 和性能计数器形成板级诊断闭环。

\begin{figure}[htb]
\centering
\includegraphics[width=\textwidth]{\detokenize{img/picture/工作流程图.png}}
\caption{LASA 系统框架与优化路线}
\label{fig:workflow_overview}
\end{figure}

\section{本文主要工作}

围绕上述问题，本文完成的主要工作如下。

\begin{enumerate}
    \item \textbf{提出并实现 LASA 层自适应脉动阵列加速器。} LASA 基于 dual-lane INT8 systolic array，支持 3x3 卷积、原生 1x1 卷积、INT32 partial-PSUM、逐通道重量化、可旁路 activation LUT、可旁路 2x2 maxpool 和 OFM 写回；系统接口采用 AXI-Lite 配置和四路 AXI-Stream/DMA 数据通路，能够与 A53 bare-metal runtime 协同执行 Conv0--Conv9。

    \item \textbf{提出 KCS 三维分块执行模型。} 本文将不同卷积层统一映射为 K-pass、Cout block 和 spatial tile 三个调度维度。K-pass 适配阵列输入 lane 数，Cout block 适配 dual-lane 输出通道并行度，spatial tile 约束片上缓存容量和 DMA 搬运粒度，为后续数据流优化提供统一执行抽象。

    \item \textbf{提出面向非计算开销的三类优化策略。} 面向 PS/PL 交互开销，BSD 通过 batch stream、packet count 和预打包权重减少 DMA 启动与软件服务次数；面向 IFM 布局转换开销，OCRR 通过 raw-HWC tile cache、1x1/3x3 replay 和 backend full-tile 降低软件重排与 tile 边界成本；面向 K-pass 边界气泡，OPF-P 通过流水化 drain、continuous PSUM collector、pass-level prefetch 和 during-compute prefetch 重叠部分和反馈与下一 pass 准备。

    \item \textbf{建立从 RTL 到 KV260 板级运行的验证和性能诊断体系。} 本文使用 Icarus Verilog、Vivado XSIM、外部 RTL golden、链式 semantic golden 和 Vitis board smoke 构成多层验证流程；在硬件中加入 stage、sub-stage、raw-HWC、prefetch、collector、pass timeline 等计数器，用于定位性能瓶颈和解释优化收益。
\end{enumerate}

\section{论文组织结构}

本文共分为六章。第一章介绍研究背景、相关研究现状和本文贡献。第二章给出卷积计算、INT8 量化、systolic array、YOLOv3-tiny 层结构和 MPSoC 平台基础。第三章提出 LASA 层自适应脉动阵列架构，介绍 PS/PL 职责划分、MPSoC 总体结构、LASA 数据通路、层自适应机制和链式推理执行流程。第四章讨论面向链式推理的数据流与调度优化，包括 KCS 三维分块执行模型、BSD 软硬件流式协同、OCRR 片上重排与重放、OPF-P 反馈重叠与预取。第五章介绍运行时集成、验证机制与性能诊断，包括 layer descriptor、STREAM\_CFG 配置、golden 验证路径、性能计数器和最终推荐软硬件组合。第六章给出正确性验证、板级端到端性能、消融实验、资源时序结果和后续工作讨论。

    \chapter{背景与相关工作}
\label{chap:related}

本章介绍与本文设计相关的基础理论与技术背景，包括卷积运算的数学定义及其 GEMM 转化、脉动阵列架构分类、INT8 量化原理、FPGA CNN 加速器相关工作以及目标硬件平台。

\section{卷积神经网络与卷积运算}

\subsection{标准卷积运算}

标准卷积层是 CNN 中最核心的计算算子。给定输入特征图（Input Feature Map, IFM） $I \in \mathbb{R}^{C_{in} \times H \times W}$、卷积核 $W \in \mathbb{R}^{C_{out} \times C_{in} \times K_h \times K_w}$ 和偏置 $B \in \mathbb{R}^{C_{out}}$，输出特征图（Output Feature Map, OFM） $O \in \mathbb{R}^{C_{out} \times H_{out} \times W_{out}}$ 的计算公式为：

\begin{equation}
O[c_{out}][y][x] = B[c_{out}] + \sum_{c_{in}=0}^{C_{in}-1} \sum_{k_y=0}^{K_h-1} \sum_{k_x=0}^{K_w-1} I[c_{in}][y \cdot S + k_y - P][x \cdot S + k_x - P] \cdot W[c_{out}][c_{in}][k_y][k_x]
\label{eq:conv2d}
\end{equation}

其中，$S$ 为步长（Stride），$P$ 为填充（Padding），$H_{out} = (H + 2P - K_h) / S + 1$，$W_{out} = (W + 2P - K_w) / S + 1$。

YOLOv3-tiny 网络中主要使用 3×3 卷积（$K_h=K_w=3$）和 1×1 卷积（通道变换层）。本文加速器专门针对 3×3 卷积的硬件友好展开进行了优化，1×1 卷积可作为 3×3 卷积的特例进行处理。

\subsection{卷积到 GEMM 的映射}

卷积运算可以通过 im2col（Image to Column）变换转化为矩阵乘法（GEMM）形式。将 IFM 中与每个滑动窗口位置对应的 $C_{in} \times K_h \times K_w$ 个元素展开为一行，将 $C_{out}$ 组卷积核分别展开为一列，即可将式~\ref{eq:conv2d}转化为：

\begin{equation}
\text{OFM}[p, c_{out}] = \text{Bias}[c_{out}] + \sum_{k=0}^{K_{total}-1} \text{IFM}[p, k] \cdot W[k, c_{out}]
\label{eq:gemm}
\end{equation}

其中，$p = oy \times W_{out} + ox$ 为输出像素的一维坐标，$K_{total} = C_{in} \times K_h \times K_w$ 为展开后的输入维度。这一映射将 6 层嵌套循环（输出通道、输出高度、输出宽度、输入通道、卷积核高度、卷积核宽度）约简为标准的矩阵乘法问题，极大便利了硬件加速器设计。

\section{脉动阵列架构}

\subsection{基本概念}

脉动阵列（Systolic Array）是一种由规整排列的处理单元（Processing Element, PE）组成的二维计算网格。数据以脉动方式在相邻 PE 间逐周期传递，每个 PE 在每个时钟周期执行固定模式的乘加运算。脉动阵列的主要优点包括：

\begin{enumerate}
    \item \textbf{规整性}：所有 PE 结构相同，适合 VLSI/FPGA 大规模阵列化实现。
    \item \textbf{数据复用}：权重、输入和部分和在阵列中以脉动传播，大幅减少对全局存储的带宽需求。
    \item \textbf{可扩展性}：通过改变阵列行列数即可调整计算吞吐与资源开销的折中。
\end{enumerate}

\subsection{数据流分类}

根据 Sze 等人~\cite{sze2017efficient}的分类体系，脉动阵列的数据流可按权重、输入和部分和的驻留策略分为以下四类：

\begin{enumerate}
    \item \textbf{权重驻留（Weight-Stationary）}：每个 PE 存储一个或多个权重值，IFM 在行方向流动，PSUM 在列方向流动。代表作为 Google TPU~\cite{jouppi2017tpu}。优点是权重从外部内存加载后驻留在 PE 中，适合权重重用率较高的场景。

    \item \textbf{输出驻留（Output-Stationary）}：每个 PE 负责累加一个输出像素的所有乘积累积，IFM 和权重在阵列中流动。优点是部分和驻留在 PE 内，避免了 PSUM 的片外存储需求。

    \item \textbf{输入驻留（Input-Stationary）}：每个 PE 存储一个 IFM 元素，权重在阵列中流动。适合 IFM 数据量较小的场景。

    \item \textbf{行静态（Row-Stationary）}：Eyeriss~\cite{chen2016eyeriss} 提出的数据流，在 1D PE 行维度保持 IFM 元素驻留，在 2D PE 列维度保持 filter 行元素驻留，以最大化本地数据复用并最小化数据搬运能量。
\end{enumerate}

本文选择权重驻留数据流，原因如下：（1）YOLOv3-tiny 的 3×3 卷积核权重总量相对有限，适合驻留在 PE 内部；（2）权重驻留方案的控制逻辑更为简洁，有利于 FPGA 实现中减少 LUT 开销；（3）在 K 分块多轮累加场景下，权重可在多个 PSUM 反馈轮次间保持驻留不变，减少了权重重载次数。

\section{INT8 量化}

神经网络的量化是指将浮点权重和激活值映射到低位宽整数的过程。Jacob 等人提出的整数量化方案~\cite{jacob2018quantization}将推理过程完全转换为整数运算：

\begin{equation}
q = \text{clamp}\left(\text{round}\left(\frac{r}{S}\right) + Z,\ q_{min},\ q_{max}\right)
\label{eq:quant}
\end{equation}

其中 $r$ 为浮点值，$S$ 为缩放因子（Scale），$Z$ 为零点（Zero Point），$q_{min}$ 和 $q_{max}$ 为量化范围（对于 INT8 为 -128 和 127）。

在推理时，卷积计算完全以整数完成：

\begin{equation}
q_{out} = Z_{out} + \frac{S_w S_{in}}{S_{out}} \sum (q_w - Z_w)(q_{in} - Z_{in})
\label{eq:int8conv}
\end{equation}

其中，预计算的乘数因子 $\frac{S_w S_{in}}{S_{out}}$ 可转化为定点乘加移位操作。本文加速器通过重量化（Requantization）模块，以每一路输出通道一个 16-bit 乘数因子加 4-bit 移位量的参数化方式实现了此量化方案。

\section{FPGA CNN 加速器相关工作}

FPGA CNN 加速器的研究可追溯至 Qiu 等人~\cite{qiu2016going}的工作，该文首次系统性地探索了在嵌入式 FPGA 上部署 CNN 的完整设计流程。随后 Guo 等人在此基础上发布了 Angel-Eye 工具链~\cite{guo2019angel}，实现了从 Caffe 模型到 FPGA 比特流的自动化映射。

在脉动阵列 FPGA 实现方面，Wei 等人~\cite{wei2017automated}提出了自动化脉动阵列综合框架，支持从高层描述生成针对不同 FPGA 器件的 RTL 代码。Ding 等人~\cite{ding2022tinyyolo}报道了一款面向 YOLOv3-tiny 的轻量级 FPGA CNN 加速器设计。然而，这些工作大多侧重于高层设计空间探索或单一模块优化，缺乏从 PE 级 DSP 原语到顶层 AXI-Stream 系统接口的端到端闭环设计和验证数据。

\section{Zynq UltraScale+ XCK26 平台}

本文目标平台为 Xilinx Kria K26 核心 FPGA 器件 XCK26（具体型号 xck26-sfvc784-2LV-c），属于 Zynq UltraScale+ MPSoC 系列。其主要逻辑资源如表~\ref{tab:xck26-resources}所示。

\begin{table}[htb]
\centering
\caption{XCK26 (sfvc784-2LV-c) 主要逻辑资源}
\label{tab:xck26-resources}
\begin{tabular}{lc}
\toprule
\textbf{资源类型} & \textbf{可用数量} \\
\midrule
系统逻辑单元 (SLC) & 约 256K \\
CLB LUT          & 117,120 \\
CLB 触发器       & 234,240 \\
Block RAM Tile   & 144 (36Kb 等效) \\
DSP48E2 Slice    & 1,248 \\
用户 I/O         & 468 \\
\midrule
内部互联总线      & AXI4-Lite / AXI4-Stream / AXI4-MM \\
PS 侧处理器       & Quad-core ARM Cortex-A53 \\
\bottomrule
\end{tabular}
\end{table}

XCK26 的 DSP48E2 是一款高性能乘累加原语，支持 27×18-bit 乘法器和 48-bit 累加器~\cite{xilinx2023dsp}。其关键特性之一是通过 (A+D)×B 的预加法运算模式，可以在单个 DSP 内完成两项独立乘法的合并计算，为本文的双 INT8 乘累加 PE 设计提供了硬件基础。

硬件平台的另一关键特性是完备的 AXI 互联体系。Zynq UltraScale+ 的处理系统（Processing System, PS）与可编程逻辑（Programmable Logic, PL）之间通过 AXI 接口实现高带宽低延迟的数据与控制交互~\cite{arm2021axi, xilinx2023axi}。本文加速器设计的 AXI-Lite 配置接口和 AXI-Stream 数据接口即为后续 PS/DMA 协同调度的直接对接边界。

    \chapter{LASA 层自适应脉动阵列架构}
\label{chap:architecture}

本文的目标不是单独实现一个卷积算子，而是在边缘 MPSoC 上完成从 DDR 图像输入、PS 侧调度、PL 侧卷积计算、OFM 回写到 YOLO decode 的完整链式推理。为此，本文提出一种层自适应脉动阵列加速器（Layer-Adaptive Systolic Accelerator, LASA）。LASA 以统一 INT8 systolic array 为计算基础，通过 layer descriptor 和 AXI-Lite 配置支持不同卷积核、输出通道分块、K pass、量化参数、激活和池化组合，从而在同一套硬件上覆盖 YOLOv3-tiny Conv0--Conv9 的多种层形状。

\section{系统设计目标与软硬件分工}

真实 YOLOv3-tiny 链式推理对加速器提出三类要求。第一，PL 侧卷积核心需要保持稳定的 INT8/INT32 计算语义，使 Conv0--Conv9 的输出能够与 golden 逐层比较。第二，PS/PL 边界需要承载真实层的数据搬运和背压，而不是依赖小规模仿真中的理想输入。第三，系统必须暴露足够的计数器和调试接口，使端到端延迟可以被分解为 DMA、IFM 重排、阵列计算、partial-PSUM 反馈、OFM 后处理和软件 decode 等阶段。

因此，本文采用“PS 负责任务组织，PL 负责规则密集计算”的协同方式。PS 侧 A53 程序维护 layer descriptor、DDR 缓冲区、DMA 启停、输出重排和 YOLO decode；PL 侧 LASA 接收寄存器配置和 AXI-Stream 数据，在片上完成卷积、重量化、激活、池化和 OFM packet 输出。第~\ref{chap:dataflow} 章讨论的 KCS 执行模型以及 BSD、OCRR、OPF-P 三类优化策略均建立在这一软硬件边界之上。

\section{MPSoC 总体结构}

图~\ref{fig:system_overview} 给出系统级结构。DDR 中保存输入图像包、量化权重、中间特征图和输出 packet；PS A53 通过 AXI-Lite 配置 LASA，并通过四路 AXI DMA 搬运 bias、weight、IFM 和 OFM 数据。LASA 以 AXI-Stream 接收三路输入流，以 AXI-Stream 输出 OFM debug packet 到 DDR，随后由软件重排并执行 YOLO decode。

\begin{figure}[htb]
\centering
\includegraphics[width=\textwidth]{\detokenize{img/picture/MPSOC架构3.png}}
\caption{MPSoC 端到端系统结构}
\label{fig:system_overview}
\end{figure}

该结构分离控制面和数据面。AXI-Lite 只承载层参数、tile 参数、量化参数、stream 配置和性能计数器；大规模张量数据全部通过 AXI DMA 搬运。这一划分使 PL 内部数据流可以持续演进，而外部软件 ABI 和 DMA 拓扑保持稳定。

\section{LASA 数据通路概览}

图~\ref{fig:accelerator_datapath} 展示 LASA 的数据通路。输入侧包含 bias loader、weight loader 和 IFM feeder；中间是 layer scheduler、systolic array 和 partial-PSUM 存储；输出侧包含 collector/drain、requant、activation、pooling 和 OFM writer。该图只描述各功能块的连接关系，具体的 KCS 分块、片上重排、PSUM 反馈和 pass 级预取在第~\ref{chap:dataflow} 章展开。

\begin{figure}[htb]
\centering
\includegraphics[width=\textwidth]{\detokenize{img/picture/LASA数据通路.png}}
\caption{LASA 数据通路概览}
\label{fig:accelerator_datapath}
\end{figure}

LASA 的最终板级配置采用 ROWS=18、COLS=8、$Cout_{tile}=16$。阵列沿 K 维接收 18 路 IFM lane；每个列位置包含两个 INT8 output-channel lane，因此 COLS=8 对应 16 个并行输出通道。两个 lane 的乘积分别累加到独立 INT32 partial sum 中，并在 final pass 后进入逐通道重量化路径。该 dual-lane INT8 column 设计提高了输出通道维度并行度，也决定了 PSUM packet、requant lane 和 OFM writer 的基本粒度。

\section{层自适应机制}

LASA 的“层自适应”体现在四个方面。第一，计算路径原生支持 3x3 与 1x1 卷积。对 3x3 卷积，IFM feeder 根据 $oy,ox,k_y,k_x,c_i$ 生成展开后的 K 维向量；对 1x1 卷积，硬件使用 native vector path，避免将 1x1 强行映射为稀疏 3x3 后产生无效 K pass。这样 Conv7 和 Conv9 等检测头可以使用更短的 $K_{total}$。

第二，后处理路径采用可配置旁路。final K pass 输出首先进入逐通道重量化模块；重量化后的 uint8 输出进入 activation LUT，可根据层描述选择 identity、ReLU/Leaky ReLU 等语义。若当前层启用 pooling，则 OFM pooling 模块执行 $2\times2$、stride=2 的 maxpool；若不启用 pooling，则数据直接旁路到 OFM byte stream。该机制使同一套硬件可以覆盖 YOLOv3-tiny 中带 pool 和不带 pool 的层。

第三，层参数由 descriptor 驱动。软件通过 layer descriptor 配置输入/输出尺寸、kernel、stride、padding、$K_{total}$、Cout blocks、spatial tile、量化参数、activation 和 pool 标志。硬件不需要为不同层重新综合，而是通过 AXI-Lite 寄存器和 stream packet 边界切换执行语义。

第四，数据流模式可组合。LASA 的基础数据通路可以与 batch stream、raw-HWC cache、partial-PSUM overlap、continuous collector 和 during-compute prefetch 等模式组合使用。对于以 1x1/3x3 标准卷积、逐层激活和池化为主的轻量化边缘 CNN，该机制具备扩展到其他顺序型卷积层的基础；对 depthwise/group convolution、residual add、concat 等结构，则需要后续扩展专门的数据路径或软件融合策略。

\section{链式推理执行流程}

一次 Conv0--Conv9 推理从 DDR 中的图像包开始。软件读取或准备输入图像，按 layer descriptor 写入当前层的尺寸、量化、tile 和 stream 配置，随后启动 OFM S2MM DMA 以及 bias、weight、IFM 三路 MM2S DMA。LASA 在当前层结束后输出带地址的 OFM packet，软件将其重排为 HWC buffer，并作为下一层输入或最终 YOLO decode 输入。

当前 OFM 输出采用 debug packet 语义：每个有效输出 byte 携带地址并写入 DDR。该格式不是最高带宽效率的连续 HWC burst，但保留了地址信息，便于板级阶段进行逐层重排、golden 比对和错误定位。第~\ref{chap:runtime} 章将进一步说明 runtime、寄存器接口、验证模式和性能计数器如何围绕该路径组织。

\section{本章小结}

本章提出并定义了 LASA 层自适应脉动阵列架构。LASA 通过 dual-lane INT8 systolic array、1x1/3x3 原生计算路径、可旁路后处理和 descriptor-driven 配置，为 YOLOv3-tiny 链式推理提供统一硬件执行基础。下一章将在 LASA 上进一步提出 KCS 三维分块执行模型，并围绕 PS/PL 交互、IFM 数据排布和 K-pass 边界气泡展开优化。

    \chapter{面向链式推理的数据流与调度优化}
\label{chap:dataflow}

LASA 提供了层级可配置的硬件执行基础，但真实 YOLOv3-tiny 链式推理的端到端延迟仍不能只由阵列峰值吞吐决定。权重加载、IFM 打包、DMA 启动、tile 边界同步、partial-PSUM 反馈、pass 准备和 OFM 输出都会形成非计算开销。因此，本文将优化目标写为
\begin{equation}
T_{noncomp}=T_{layer}-T_{compute\_fire}.
\end{equation}
其中 $T_{compute\_fire}$ 表示阵列真正接受输出 pixel 计算的周期数，其余部分均属于等待、搬运、反馈或边界开销。

本章首先提出 KCS 三维分块执行模型，将不同层统一映射为 K-pass、Cout block 和 spatial tile。随后围绕三类非计算瓶颈提出三类优化策略：面向 PS/PL 交互的批量流式协同策略 BSD（Batched Streaming Dataflow）、面向 IFM 布局转换的片上重排与重放策略 OCRR（On-Chip Reorder and Replay），以及面向 K-pass 边界气泡的反馈重叠与预取策略 OPF-P（Overlapped PSUM Feedback and Prefetch）。

\begin{figure}[htb]
\centering
\includegraphics[width=\textwidth]{\detokenize{img/picture/时序优化策略.png}}
\caption{BSD、OCRR 与 OPF-P 三类优化策略总览}
\label{fig:strategy_overview}
\end{figure}

\section{KCS 三维分块执行模型}

KCS 执行模型由 K 维、输出通道和空间三个调度维度组成。K 维分块由 ROWS 决定，最终配置 ROWS=18，因此 $K_{tile}=18$。输出通道分块由 COLS 和 dual-lane INT8 column 决定，每列产生两个输出通道 lane，因此
\begin{equation}
Cout_{tile}=2\times COLS.
\end{equation}
最终配置 COLS=8，对应 $Cout_{tile}=16$。空间分块由 tile 的输出行范围决定，用于限制片上 IFM/OFM buffer 和单次 AXI DMA 长度。

图~\ref{fig:tiling_schedule} 给出三维调度关系。每个 spatial tile 内，调度器先遍历输出通道块，再对当前输出通道块执行所有 K pass。首个 K pass 从 bias 初始化 partial sum；非首 pass 读取上一轮 partial-PSUM；final pass 将结果送入 requant、activation、pooling 和 OFM path。

\begin{figure}[htb]
\centering
\includegraphics[width=\textwidth]{\detokenize{img/picture/KCS分块策略.png}}
\caption{KCS 三维分块执行模型}
\label{fig:tiling_schedule}
\end{figure}

KCS 不是单独的优化策略，而是 LASA 上的统一映射抽象。它回答不同 CNN 层如何映射到固定阵列规模，而 BSD、OCRR 和 OPF-P 则分别针对映射后暴露出的 PS/PL 交互、数据排布和 pass 边界开销进行优化。

\section{BSD：面向 PS/PL 交互的批量流式协同}

BSD 的目标是降低 PS 逐请求服务和 DMA 频繁启动造成的系统协同开销。早期 per-request 模式中，硬件每需要一个 bias block、weight tile 或 IFM line 时向软件发出请求，软件随后打包数据、启动 DMA 并等待完成。这种方式适合小型 smoke test，但在 Conv6 等真实层中会产生大量服务事件，A53 侧开销和 DMA 启动开销会淹没阵列计算。

BSD 将每个 spatial tile 所需的 bias、weight 和 IFM 数据预先打包成连续 AXI stream。软件在 tile 开始时分别启动三路输入 DMA，硬件通过 packet count 恢复内部边界，只在整个 stream 末尾依赖 TLAST。这样，DMA 启动次数从每个请求一次降低为每个 tile 每类输入一次，AXI backpressure 也由 PL 自然吸收。

权重数据在模型中通常按输出通道、输入通道和 kernel 维度组织，而硬件消费顺序是 $Cout block \rightarrow K pass \rightarrow lane \rightarrow channel$。本文在生成 Vitis header 时直接产生硬件消费顺序的权重数组，使 weight DMA 可以读取 ELF 只读段中的连续数组，不再需要运行时 scratch copy 或 repack。PL 端 weight loader 使用 64-bit AXI beat 写入多个 byte bank，进一步降低 weight load 等待。预打包权重虽然作用于数据布局，但其主要收益是减少运行时 PS 服务和 DMA 准备成本，因此归入 BSD。

\section{OCRR：面向 IFM 布局转换的片上重排与重放}

OCRR 的目标是缓解软件自然 HWC 布局与 systolic array K-lane 消费顺序之间的不匹配。IFM 输入存在两种路径。第一种是 prepacked IFM stream：软件提前把每个 K pass 需要的 IFM lane 按硬件顺序打包。该方式 PL 简单，但 A53 侧打包成本较高。第二种是 raw-HWC stream：软件发送自然 HWC tile，PL 在片上缓存后按 pass 和窗口顺序 replay。

图~\ref{fig:strategy_overview} 中间部分给出 OCRR 的数据排布思路。prepacked 路径把重排压力放在软件；raw-HWC 路径把重排压力转移到 PL cache/replay，但可减少 A53 逐 pass 复制和索引开销。对 3x3 层，cache materializes 每个输出 pixel 的 3x3 窗口，并按 18 lane 输出给阵列；对 1x1 层，cache 按连续通道向量 replay。

OCRR 还包括后端 full-tile 调度。Conv5、Conv6 和 Conv8 的空间尺寸均为 $13\times13$。早期调度将这些层拆成 $4+4+4+1$ 行 spatial tile，可以降低片上缓存容量需求，但每个 tile 都需要重复启动 bias、weight、IFM 和 OFM 流程。随着 raw-HWC cache 使用 URAM 扩容，后端层可以使用完整 $13\times13$ tile，从而减少 tile 边界固定成本。容量条件可近似写为
\begin{equation}
tile\_pixels \times \left\lceil\frac{C_{in}}{2}\right\rceil \le HWC\_CACHE\_DEPTH.
\end{equation}
在最终 full-tile build 中，Conv5/8 需要 $13\times13\times\lceil256/2\rceil=21632$ 个 materialized word，Conv6 需要 $13\times13\times\lceil512/2\rceil=43264$ 个 materialized word。raw-HWC 并非对所有层都收益为正；Conv3 raw-HWC 虽然可功能通过，但由于 tile/FIFO 容量和 replay 开销，板级结果不优于推荐配置，因此作为边界分析而非默认路径。

\section{OPF-P：面向 K-pass 边界的反馈重叠与预取}

OPF-P 的目标是隐藏多 K pass 之间的 partial-PSUM 反馈和下一 pass 准备气泡。对于 $K_{total}>K_{tile}$ 的层，第一轮 K pass 只能产生部分和，后续 pass 必须读取前一轮 partial sum 并继续累加，最后一轮才进入重量化和输出路径。对一个输出 pixel 和一个 $Cout_{tile}$ 通道块，硬件内部形成
\begin{equation}
P[p,c] = \sum_{k=k_0}^{k_0+K_{tile}-1} X[p,k]\cdot W[k,c].
\end{equation}
当当前 pass 不是 final pass 时，该 packet 被写入片上 partial-PSUM buffer；下一 pass 读取同一地址的 partial sum 作为初值继续累加。当当前 pass 是 final pass 时，该 packet 被送入 requant/activation/pool/OFM path。

该语义要求 partial-PSUM 地址、Cout block 和 pass context 与下一 pass 严格对齐，同时避免读写冲突，并保持 final pass 输出顺序与软件重排和 golden 比对一致。本文使用 ping-pong bank 和 per-pass context 维护这些条件。也就是说，PSUM 反馈不是独立于数据流之外的模块，而是跨 K pass 数据流的一部分。

若每轮 pass 都严格串行执行“计算完成、全部 drain、下一轮 feed、下一轮 compute”，后端大层会出现明显阵列空泡。本文首先将 drain writer 改为 AXI-Stream 风格握手：packet valid、addr 和 data 在真实 $valid \& ready$ 握手前保持稳定，握手成功后才推进地址。随后加入 read-return tracker 和 one-entry hold register，使下游 ready 时可以接近每周期输出一个 PSUM packet。

Early drain 允许调度器在当前 pass 已经开始产生 PSUM 数据后提前启动 drain，只要 partial-PSUM FIFO 已具备可读数据。Continuous PSUM collector 进一步改变非 final pass 的路径：collector 在 compute 开始前接收 pass context，并在阵列列 FIFO 产生完整 packet 后立即将其写入 partial-PSUM bank；final pass packet 则进入后处理流水线。这样，非 final pass 的 partial-PSUM 写回从“compute 后集中 drain”变为“compute 过程中持续收集”，减少 pass 边界处的反馈延迟。

K pass 之间除了 PSUM 反馈，还包含下一 pass 权重加载和 IFM replay。Pass-level prefetch 在当前 pass 已经满足安全条件后提前准备下一 pass 的 weight 和 raw-HWC replay。该机制遵循两个约束：一是不能覆盖当前 pass 正在使用的 PE 权重；二是 next-pass IFM replay 只能进入解耦 FIFO 或安全队列，不能破坏当前 pass 的输入顺序。During-compute prefetch 进一步允许在当前 pass compute 进行时启动下一 pass 权重和 raw-HWC replay 的 staging，从而降低 compute stage 中由于下一 pass 准备不足造成的 idle。

图~\ref{fig:psum_prefetch_timeline} 对比了串行 pass 和 OPF-P 重叠 pass 的时间线。

\begin{figure}[htb]
\centering
\includegraphics[width=\textwidth]{\detokenize{img/picture/partial-psum feedback.png}}
\caption{OPF-P 反馈重叠与 pass 级预取时间线}
\label{fig:psum_prefetch_timeline}
\end{figure}

\section{优化策略小结}

表~\ref{tab:strategy_summary} 汇总了 KCS 执行模型与三类优化策略的定位。KCS 负责将不同卷积层统一映射到 LASA；BSD 作用于软硬件协同层，减少 PS/DMA 服务开销；OCRR 作用于数据排布层，缓解 IFM 布局转换和 tile 边界开销；OPF-P 作用于流水线调度层，隐藏 K-pass 反馈和下一 pass 准备气泡。

\begin{table}[htb]
\centering
\caption{KCS 执行模型与三类优化策略}
\label{tab:strategy_summary}
\begin{tabular}{@{}p{0.12\textwidth}p{0.16\textwidth}p{0.26\textwidth}p{0.30\textwidth}@{}}
\toprule
名称 & 层次 & 主要瓶颈 & 对应机制 \\
\midrule
KCS & 执行模型 & 层形状到固定阵列的映射 & K-pass、Cout block、spatial tile \\
BSD & 软硬件协同 & PS 服务与 DMA 启动 & batch stream、packet count、prepacked weight \\
OCRR & 数据排布转换 & IFM 重排与 tile 边界 & raw-HWC cache、replay、backend full-tile \\
OPF-P & 流水线调度 & PSUM/pass 边界气泡 & drainpipe、continuous collector、pass prefetch \\
\bottomrule
\end{tabular}
\end{table}

三类优化共同目标是减少 $T_{noncomp}$，但它们针对的不是同一种开销。这样的分层表述使本文的优化链条从零散实现点收敛为系统协同、数据排布和流水线调度三个层面的策略。下一章将说明这些策略如何通过 runtime、寄存器接口、验证流程和性能计数器集成为可运行、可诊断的板级系统。

    \chapter{Partial-PSUM 反馈与系统集成}
\label{chap:psum}

多 K pass 卷积的核心难点在于 partial-PSUM feedback。对于 $K_{total}>K_{tile}$ 的层，第一轮 K pass 只能产生部分和，后续 pass 必须读取前一轮 partial sum 并继续累加，最后一轮才进入重量化和输出路径。若每轮 pass 都严格串行执行“计算完成、全部 drain、下一轮 feed、下一轮 compute”，后端大层会出现明显阵列空泡。本章介绍本文如何对 partial-PSUM 路径和 pass 边界进行流水化与重叠。

\section{Partial-PSUM 反馈语义}

对一个输出 pixel 和一个 $Cout_{tile}$ 通道块，硬件内部形成一个 PSUM packet：
\begin{equation}
P[p,c] = \sum_{k=k_0}^{k_0+K_{tile}-1} X[p,k]\cdot W[k,c].
\end{equation}
当当前 pass 不是 final pass 时，该 packet 被写入片上 partial-PSUM buffer；下一 pass 读取同一地址的 partial sum 作为初值继续累加。当当前 pass 是 final pass 时，该 packet 被送入 requant/activation/pool/OFM path。

该语义要求三个条件同时成立：第一，partial-PSUM 的地址和 Cout block 必须与下一 pass 对齐；第二，写入和读取不能发生读写冲突；第三，final pass 的输出顺序必须与软件重排和 golden 比对一致。本文使用 ping-pong bank 和 per-pass context 维护这些条件。

\section{流水化 PSUM Drain}

早期 drain writer 在下游 ready 出现回压时容易产生地址推进和 valid 保持之间的时序风险，同时吞吐受内部状态机限制。本文将 drain writer 改为 AXI-Stream 风格握手：packet valid、addr 和 data 在真实 $valid \& ready$ 握手前保持稳定，握手成功后才推进地址。随后加入一拍 read-return tracker 和 one-entry hold register，使下游 ready 时可以接近每周期输出一个 PSUM packet。

流水化 drain 的意义不仅是提高 drain 吞吐，也为后续 overlap 提供可控事件：调度器可以基于 drain packet fire、drain done 和 FIFO empty wait 判断 partial-PSUM 是否已经产生足够可读数据。

\section{Early Drain 与 Continuous PSUM Collector}

传统 pass 调度在 compute done 后才启动 drain。Early drain 允许调度器在当前 pass 已经开始产生 PSUM 数据后提前启动 drain，只要 partial-PSUM FIFO 已经具备可读数据。该机制降低 compute 完成后等待 drain 启动的空泡，但仍保留 final pass 和非 final pass 的输出顺序约束。

Continuous PSUM collector 进一步改变非 final pass 的路径。collector 在 compute 开始前接收一个 pass context，并在阵列列 FIFO 产生完整 packet 后立即将其写入 partial-PSUM bank；final pass packet 则进入后处理流水线。这样，非 final pass 的 partial-PSUM 写回从“compute 后集中 drain”变为“compute 过程中持续收集”，减少 pass 边界处的反馈延迟。

图~\ref{fig:psum_prefetch_timeline} 对比了串行 pass 和重叠 pass 的时间线。

\begin{figure}[htb]
\centering
\resizebox{\textwidth}{!}{%
\begin{tikzpicture}[
    x=8mm,y=7mm,
    phase/.style={draw, minimum height=5mm, align=center},
    arr/.style={-Latex, thick}
]
\node at (0,2) {串行};
\node[phase, minimum width=14mm] at (2,2) {feed};
\node[phase, minimum width=18mm] at (4.1,2) {compute};
\node[phase, minimum width=18mm] at (6.4,2) {drain};
\node[phase, minimum width=16mm] at (8.5,2) {next feed};
\node[phase, minimum width=20mm] at (10.8,2) {next compute};

\node at (0,0) {重叠};
\node[phase, minimum width=14mm] at (2,0) {feed};
\node[phase, minimum width=28mm] at (4.6,0) {compute};
\node[phase, minimum width=24mm] at (5.3,-0.8) {continuous collect};
\node[phase, minimum width=24mm] at (7.6,0.8) {next prefetch};
\node[phase, minimum width=22mm] at (9.5,0) {next compute};
\draw[arr] (6.9,0.8) -- (8.5,0.1);
\end{tikzpicture}
}
\caption{Partial-PSUM 反馈与 pass-level prefetch 时间线}
\label{fig:psum_prefetch_timeline}
\end{figure}

\section{Pass-Level Prefetch}

K pass 之间除了 PSUM 反馈，还包含下一 pass 权重加载和 IFM replay。Pass-level prefetch 在当前 pass 已经满足安全条件后提前准备下一 pass 的 weight 和 raw-HWC replay。该机制遵循两个约束：一是不能覆盖当前 pass 正在使用的 PE 权重；二是 next-pass IFM replay 只能进入解耦 FIFO 或安全队列，不能破坏当前 pass 的输入顺序。

本文实现了两级 prefetch。第一阶段 pass prefetch 主要在当前 compute 完成后、drain 或 pass 边界窗口中准备下一 pass；第二阶段 during-compute prefetch 进一步允许在当前 pass compute 进行时启动下一 pass 权重和 raw-HWC replay 的 staging。后者能够降低 compute stage 中由于下一 pass 准备不足造成的 idle，但需要更严格地维护权重读取预算和 IFM FIFO 顺序。

\section{STREAM\_CFG 与调度开关}

所有优化均通过 AXI-Lite 中的 STREAM\_CFG 位控制。bit0 使能 batch mode，bit1 使能 raw-HWC mode，bit2 使能 early drain，bit3 使能 pass prefetch，bit4 使能 partial-PSUM overlap，bit5 使能 continuous PSUM collector，bit7 使能 during-compute prefetch。bit6 对应 column-level PSUM，是当前探索性设计，不作为最终推荐配置的必要条件。

这种设计保持了软件 ABI 的稳定性。默认关闭实验开关时，系统回到保守串行语义；逐步打开开关时，软件只需写入不同 STREAM\_CFG 和少量 tail/prefetch 参数，外部 AXI packet 格式保持不变。该策略便于在同一 runtime 中进行 A/B 测试和板级回归。

\section{性能计数器设计}

为了避免仅凭端到端延迟判断优化效果，本文在硬件中加入多层性能计数器。

\begin{table}[htb]
\centering
\caption{主要性能计数器类别}
\label{tab:perf_counters}
\begin{tabular}{ll}
\toprule
类别 & 作用 \\
\midrule
PERF & 统计 tile busy、外部等待、compute fire 等顶层周期 \\
STAGEPERF & 统计 bias、weight、feeder、compute、drain、OFM post 阶段 \\
SUBPERF & 细分 feeder push、FIFO stall、compute active、IFM stall、tail 等 \\
RAWSTAT & 统计 raw-HWC load、unpack、replay 和 replay wait \\
PREFETCHPERF & 统计 prefetch start、hit、miss、stall 等事件 \\
COLLECTPERF & 统计 collector packet、partial write、final write 和 context stall \\
PASSPERF & 统计 pass start 到 first fire、fire span、compute idle 等 \\
\bottomrule
\end{tabular}
\end{table}

这些计数器使本文能够说明每次优化后的瓶颈迁移。例如，drainpipe 主要降低 drain 阶段；continuous collector 降低 partial-PSUM 存储开销；backend full-tile 减少 spatial tile 固定成本；during-compute prefetch 降低下一 pass 准备造成的 compute idle。第六章将基于这些计数器进行实验分析。

\section{探索性 Column-Level PSUM}

当前 continuous PSUM collector 仍以完整 packet 为反馈粒度：collector 等待所有列 FIFO 都有数据后写入一个 $COLS\times2$ 通道 packet。Column-level PSUM 的探索方向是按列独立写入和读取 partial sum，使下一 pass 能够按照 systolic array 的列传播延迟更早获得可用数据。

本文已经实现了 column-level ping-pong buffer 和 column stream feeder 的基础模块，并完成模块级验证。但该路径尚未作为最终推荐板级配置，因此本文将其作为未来工作讨论，不将其计入核心性能贡献。这样可以避免把尚未完整端到端验证的实验性设计包装为成熟成果。

\section{本章小结}

本章从 partial-PSUM 反馈角度解释了后端大层的性能瓶颈，并介绍了 drainpipe、early drain、continuous collector、pass prefetch 和 during-compute prefetch 等调度优化。与第四章的数据搬运优化结合后，系统能够同时减少软件打包、DMA 启动、tile 边界、PSUM drain 和 pass 准备空泡，为第六章的端到端性能提升奠定基础。

    \chapter{验证与实验结果}
\label{chap:results}

本章报告本文加速器的正确性验证、板级端到端性能、优化消融、阶段计数器分析和资源时序结果。实验的重点不是将最终延迟包装为实时视频领先指标，而是验证本文提出的数据流和调度优化是否能够在真实 MPSoC 环境中降低 YOLOv3-tiny 链式推理的非计算开销，并解释性能瓶颈如何迁移。

\section{实验设置}

实验平台为 Xilinx Kria KV260，PL 器件为 XCK26，工具链为 Vivado/Vitis 2022.2。最终推荐硬件构建目录为 \texttt{D:/MPSoC/b\_kprefetch\_22}，核心时钟目标为 100 MHz。加速器参数为 ROWS=18、COLS=8、$Cout_{tile}=16$、IFM FIFO depth=1024，后端 raw-HWC full-tile cache 使用 URAM 承载 Conv5/6/8 的 $13\times13$ tile。

最终推荐软件配置为 \texttt{conv0\_conv9\_ddr\_demo}，开启 RawHwcConv4/5/6/8、early drain、pass prefetch、partial-PSUM overlap、continuous PSUM、backend full-tile 和 during-compute prefetch。验证路径包含两类：第一类是 batch-chain 模式，对 Conv0--Conv9 每层输出与 RTL semantic golden 做 bit-exact 比较；第二类是 DDR image demo，将 $416\times416\times3$ 图像包写入 DDR，由硬件执行十层推理并在 A53 侧进行 YOLO decode。

\section{正确性验证}

本文采用从模块到板级的多层验证。RTL 层面使用 Icarus Verilog 和 Vivado XSIM 覆盖 PE、array、feeder、AXI loader、PSUM buffer、collector、AXI-Lite bridge 和端到端 conv top。系统层面使用外部 golden 数据验证真实 Conv0、Layer06/Conv3、Conv4 以及后端 Conv5/6/8 raw-HWC tile。板级层面使用 KV260 batch-chain 和 DDR image demo 验证完整链路。

\begin{table}[htb]
\centering
\caption{关键正确性验证结果}
\label{tab:correctness}
\begin{tabular}{p{0.23\textwidth}p{0.49\textwidth}p{0.18\textwidth}}
\toprule
验证对象 & 覆盖内容 & 结果 \\
\midrule
Conv0 crop + pool & 真实输入 crop、卷积、LUT、pool、OFM packet & 0 mismatch \\
Layer06/conv3 pool & $52\times52\times64 \rightarrow 52\times52\times128$ & XSIM 与板级通过 \\
Conv4 pool & $26\times26\times128 \rightarrow 13\times13\times256$ & XSIM 与板级通过 \\
Conv5/6/8 raw-HWC & 后端 full-tile raw-HWC + overlap 配置 & XSIM 2726/0 \\
Conv0--Conv9 batch-chain & 十层链式 RTL semantic golden & PASS \\
DDR image demo & 动态图像输入与 YOLO decode & 两张图像 PASS \\
\bottomrule
\end{tabular}
\end{table}

需要强调的是，链式 golden 与单层 standalone golden 不可混用。由于前级硬件量化、激活和池化语义会影响后级输入，Conv5 之后的 expected output 必须由同一条 RTL semantic chain 生成。本文在 batch-chain 中逐层保存硬件输出并比较，从而验证层间数据重排、feature buffer 切换和检测头输出的一致性。

\section{端到端性能与优化消融}

表~\ref{tab:ablation} 汇总了主要板级优化版本的端到端 DDR demo 延迟。所有结果均来自 KV260 板级运行，输入图像和模型链路保持一致。不同版本的目标不是形成严格产品发布序列，而是展示本文优化如何逐步降低非计算开销。

\begin{table}[htb]
\centering
\caption{端到端 DDR demo 优化消融}
\label{tab:ablation}
\begin{tabular}{p{0.27\textwidth}p{0.43\textwidth}p{0.18\textwidth}}
\toprule
版本或优化阶段 & 主要变化 & 典型延迟 \\
\midrule
Batch/DDR baseline & batch stream，DDR 图像 demo 初版 & $\approx 1.18$ s \\
Weight 64-bit loader & 预打包权重与 64-bit PL 解包 & $\approx 0.861$ s \\
Drainpipe & 流水化 PSUM drain writer & $\approx 0.646$ s \\
Early drain + pass prefetch & 提前 drain 与下一 K pass 预取 & $\approx 0.387$ s \\
Continuous PSUM collector & 非 final pass 持续写 partial-PSUM & $\approx 0.371$ s \\
Backend full-tile & Conv5/6/8 使用 $13\times13$ full tile & $\approx 0.335$ s \\
During-compute prefetch & 当前 compute 中准备下一 K pass & $\approx 0.288$ s \\
\bottomrule
\end{tabular}
\end{table}

最终 \texttt{b\_kprefetch\_22} 构建完成 full-programming batch-chain 验证，日志中 UART detections 与 decode golden 匹配。两张 DDR image demo 的结果如表~\ref{tab:ddr_demo} 所示。

\begin{table}[htb]
\centering
\caption{最终 DDR image demo 结果}
\label{tab:ddr_demo}
\begin{tabular}{llll}
\toprule
图像 & 总延迟 & 检测结果 & 日志 \\
\midrule
maksssksksss0.png & 288.002 ms & with\_mask, score=0.357321 & 20260616\_193623 \\
maksssksksss1.png & 287.993 ms & with\_mask, score=0.295050 & 20260616\_193752 \\
\bottomrule
\end{tabular}
\end{table}

该结果约为 3.5 FPS，不应表述为通用实时检测领先性能。本文将其解释为在资源受限 XCK26 上，通过数据流与调度优化将早期约 1.18 s 的完整链式推理降低到约 0.29 s，获得约 4.1 倍端到端加速。

\section{阶段计数器分析}

仅报告端到端延迟无法说明优化来源。表~\ref{tab:stage_breakdown} 给出最终 DDR demo 的关键硬件计数器换算时间。可以看到，最终硬件 busy 时间约为 247.184 ms，端到端 total 与 PL busy 之间的差额来自 A53 runtime、DMA 管理、cache maintenance、OFM 重排和 YOLO decode 等软件侧开销。

\begin{table}[htb]
\centering
\caption{最终版本关键阶段计数器}
\label{tab:stage_breakdown}
\begin{tabular}{lr}
\toprule
指标 & 时间 \\
\midrule
Total DDR demo & 288.002 ms \\
HW busy & 247.184 ms \\
compute\_fire & 74.323 ms \\
feeder stage & 99.913 ms \\
compute stage & 121.022 ms \\
compute idle in stage & 46.699 ms \\
drain stage & 16.467 ms \\
\bottomrule
\end{tabular}
\end{table}

这些计数器说明，经过 drainpipe 和 continuous PSUM 后，drain 已不再是主要瓶颈；经过 backend full-tile 和 during-compute prefetch 后，剩余开销主要分布在 feeder/replay 和 compute stage idle。换言之，下一步优化不应继续单纯扩大阵列，而应继续提高 raw-HWC replay 吞吐、减少 compute stage 内部空泡，并评估 OFM debug packet 对软件重排的影响。

\section{资源与时序}

表~\ref{tab:resource} 给出最终推荐硬件构建的实现资源与时序。该版本在 Vivado 2022.2 下完成布局布线，WNS 为正，说明 100 MHz 目标可闭合。

\begin{table}[htb]
\centering
\caption{最终推荐构建资源与时序}
\label{tab:resource}
\begin{tabular}{lr}
\toprule
项目 & 数值 \\
\midrule
Build directory & \texttt{D:/MPSoC/b\_kprefetch\_22} \\
CLB LUTs & 84480 \\
CLB Registers & 53260 \\
BRAM Tile & 63 \\
URAM & 24 \\
DSP & 184 \\
WNS & +0.092 ns \\
TNS & 0 ns \\
WHS & +0.010 ns \\
Routing errors & 0 \\
\bottomrule
\end{tabular}
\end{table}

资源结果体现了本文优化的代价。raw-HWC full-tile cache 引入 URAM 以减少后端 tile 启动开销；continuous PSUM collector 和性能计数器增加 LUT/FF；during-compute prefetch 增加调度与队列控制逻辑。整体上，设计仍能装入 XCK26 并闭合 100 MHz，但 LUT 使用已经较高，继续增加 column-level PSUM 或更深 overlap 时需要谨慎评估资源和时序裕量。

\section{与相关工作的定性比较}

表~\ref{tab:qualitative_compare} 从研究重点角度比较本文与代表性工作。由于不同论文使用的模型、输入尺寸、器件、频率、量化策略和是否计入软件后处理差异较大，本文不进行不公平的绝对 FPS 横向排名。

\begin{table}[htb]
\centering
\caption{与代表性工作的定性比较}
\label{tab:qualitative_compare}
\begin{tabular}{p{0.18\textwidth}p{0.24\textwidth}p{0.22\textwidth}p{0.25\textwidth}}
\toprule
特征 & Angel-Eye~\cite{guo2019angeleye} & S2A~\cite{wei2017s2a} & 本文 \\
\midrule
主要目标 & 自动化映射流程 & systolic array 生成 & MPSoC 链式推理优化 \\
网络粒度 & 多模型映射 & 阵列综合 & YOLOv3-tiny Conv0--Conv9 \\
数据流讨论 & 编译映射为主 & 阵列结构为主 & DMA/IFM/PSUM/pass 调度 \\
板级端到端 & 有 & 非重点 & KV260 DDR image demo \\
性能解释 & 较少 & 较少 & stage/pass/counter 诊断 \\
\bottomrule
\end{tabular}
\end{table}

本文的优势在于系统闭环和瓶颈解释；不足在于通用性和绝对吞吐仍有限。当前 runtime 针对 YOLOv3-tiny 单尺度静态调度，不是通用 CNN 编译器；OFM 输出仍采用 debug packet 格式，验证友好但带宽效率不足；功耗数据尚未形成板级实测，后续可结合 Vivado power estimate 或外部功耗计进行能效评估。

\section{局限性与后续工作}

本文仍存在以下局限。第一，最终延迟约 288 ms，不能宣称实时视频级性能领先。其价值在于证明优化路径和系统完整性，而不是与 GPU/NPU 进行绝对性能竞争。第二，OFM debug packet 以每字节一个带地址 packet 的方式输出，便于验证但降低 DMA 带宽效率；后续应设计连续 HWC burst 写回。第三，raw-HWC cache 对层形状敏感，Conv3 实验说明并非所有层都适合启用该路径。第四，column-level PSUM 仍处于探索阶段，需要完整 top-level xsim 和板级 A/B 后才能作为正式优化。

后续工作可从四个方向展开：一是将 OFM 输出迁移为连续 HWC DMA burst，减少软件重排；二是继续优化 raw-HWC replay 吞吐和 compute stage idle；三是完善 column-level PSUM feedback，使 partial sum 按列粒度更早进入下一 pass；四是补充功耗实测和更广泛模型验证，提升论文结论的可比性和可推广性。

\section{本章小结}

本章给出了本文加速器的正确性、性能、资源和限制分析。实验表明，本文设计能够在 KV260/XCK26 上完成真实 Conv0--Conv9 链式推理，并通过 batch stream、raw-HWC cache、backend full-tile、continuous PSUM 和 during-compute prefetch 将端到端 DDR demo 延迟降低到约 0.29 s。性能计数器进一步说明，优化后瓶颈已从 PSUM drain 和权重加载转移到 feeder/replay 与 compute idle，为后续工作提供了明确方向。


    \appendix
    \printbibliography[heading = bibintoc]

\end{document}
